% Incremental search: walk a fixed step across the search range and report
% every subinterval [x_l, x_u] where f changes sign (a bracket for a root).
% Cropped to its own bounding box so it can be dropped onto a Xournal++ page
% as an image.
%
%   latexmk -pdf incremental-search-pseudocode.tex
\documentclass[border=6pt,varwidth=11cm]{standalone}

\usepackage{algorithm2e}
\usepackage{amsmath}

\RestyleAlgo{plain}          % no float, no ruled box
\DontPrintSemicolon
\SetKwComment{Comment}{// }{}
\SetKwInOut{Input}{input}
\SetKwInOut{Output}{output}

\begin{document}
\begin{minipage}{\linewidth}
\begin{algorithm}[H]
  \textbf{Incremental Search}
  \BlankLine
  \Input{function $f$, search range $[x_\text{min}, x_\text{max}]$,
         number of steps $n$}
  \Output{list of brackets $[x_l, x_u]$, each containing a root}
  \BlankLine
  $\Delta x \leftarrow (x_\text{max} - x_\text{min}) \,/\, n$\;
  $\texttt{brackets} \leftarrow \{\,\}$\;
  $x_l \leftarrow x_\text{min}$\;
  \For{$k \leftarrow 1$ \KwTo $n$}{
    $x_u \leftarrow x_l + \Delta x$\;
    \If{$f(x_l) \cdot f(x_u) < 0$}{
      \tcp{sign change $\Rightarrow$ a root lies in $[x_l, x_u]$}
      append $[x_l, x_u]$ to \texttt{brackets}\;
    }
    $x_l \leftarrow x_u$\;
  }
  \Return \texttt{brackets}\;
\end{algorithm}
\end{minipage}
\end{document}
